\begin{description}
   \item[numeric] Refers to numbers. They can be operated on using arithmetic operations such as addition, subtraction, multiplication, and division. In terms of measurement scales, they correspond to either interval or ratio scales.
   \item[factor] Refers to the factor type. It is also known as the categorical type. It refers to the "factor" in factorial designs. This variable consists of labels that are not numerical in nature. In terms of measurement scales, it corresponds to the nominal scale.
   \item[orderd factor] Refers to the ordered factor type. It is a special case of the factor type, where the categories have a specific order. In terms of measurement scales, it corresponds to the ordinal scale. However, it is not used very frequently.
   \item[chr] Refers to the character type. At first glance, it may seem difficult to distinguish from the factor type, but this type has no corresponding numerical values and consists only of characters. If a variable is enclosed in single or double quotation marks, it will become of this type.
   \item[logical] Refers to the logical type, which can have only one of two values: \textit{TRUE} or \textit{FALSE}. \textit{TRUE} corresponds to 1, and \textit{FALSE} corresponds to 0. It is a special type of nominal scale that can be used for on/off switches and other similar applications. The labels for this type, including uppercase \texttt{TRUE} or \texttt{FALSE} and the first letters of both, \texttt{T} or \texttt{F}, are reserved words for this type, so it is better not to use them for object names, etc. Other reserved words include \texttt{NULL}, which indicates nothing; \texttt{Inf}, which represents infinity; \texttt{NA}, which represents missing values; and \texttt{pi}, a mathematical constant.
